\begingroup
\let\clearpage\relax
\let\cleardoublepage\relax
\chapter*{The argument of Part IV}
\endgroup
\addcontentsline{toc}{chapter}{The argument of Part IV}
\label{ch:results-map}

The two results chapters run a loop between solver and physics.

We first show that the trained states agree closely with the external energy references we
compare them against---an agreement rather than an independent validation, and one that stops
at $\omega=0.1$ for $N\ge6$, below which the energies are predictions. That
agreement is the precondition for using the states as instruments in the physics analysis
that follows. Read for the physics they contain, they show a crossover from compact dot to
finite Wigner molecule that is \emph{staged}: the cloud expands and its shells stiffen
radially well before the electrons acquire angular order within them.

We then turn from the state to the solver. The implemented model families have different
variational geometries even where their total energies agree. Under the parametrisations
used here, the message-passing correlator concentrates its tangent weight on fewer directions
that are better described by the selected operator dictionary. The message-passing backflow
also avoids the rank collapse seen in the smaller conventional implementation. Architecture,
optimiser and sampling measure then meet
in a single object, the geometry of the wavefunction's log-derivatives, and the evidence
suggests that natural-gradient preconditioning earns its cost only while that geometry can
still be estimated from a batch.

Finally, a generic collocation proposal loses the Wigner state. The structural
measurements of Chapter~\ref{ch:results} supply the remedy: a proposal built from the
observed shell geometry carries a cascade with MCMC-free updates to a stable state at
$\omega=0.0035$, where transfer to the next confinement fails. The component decomposition
is consistent with a disturbed Slater--Jastrow balance but does not isolate it.

\medskip
\noindent Three differently sized sets of models appear below---the production states that
carry the physics, the capacity ladder of the architecture study, and the smaller collocation
networks. They are kept distinct, and any cross-set comparison is identified as unmatched.
Table~\ref{tab:hyperparams} gives all three. The four research questions of
\S\ref{sec:research-questions} are answered here:

\begin{center}\small
\begin{tabular}{@{}ll@{}}
\toprule
\textbf{RQ4 --- Physics} & \S\ref{sec:wigner-molecule} \\
\textbf{RQ1 --- Architecture} & \S\ref{sec:kernel-q1} \\
\textbf{RQ2 --- Optimiser} & \S\ref{sec:kernel-cond} \\
\textbf{RQ3 --- MCMC-free gradient updates} & \S\ref{sec:kernel-mcmcfree}--\S\ref{sec:kernel-q3-frontier} \\
\bottomrule
\end{tabular}
\end{center}

\clearpage
